\section{Traditional and Profane Science}

\begin{quotex}
Time is the moving image of eternity.
\flright{\textsc{Plato}, \textit{Timaeus}}
\end{quotex}


Metaphysics begins where physics, that is, the study of nature, ends. We pointed out where physics, or science ends\footnote{\url{https://gornahoor.net/?p=1303\%22}}, so now let us draw the metaphysical conclusions.

\paragraph{Energy}

Taken to its extreme, in Schroedinger's model, matter consists ultimately of undifferentiated energy which takes on a concrete forms only when observed. This is strikingly similar to the notion of
\emph{Prakriti}, which Guenon describes as ``undifferentiated primordial substance'' in \emph{Man and his Becoming}. \emph{Prakriti}, as undifferentiated, comprises all forms in potential, yet cannot in itself produce or cause anything; metaphysically, then, \emph{Purusha} (spirit, consciousness) is what produces form. While Quantum Theory cannot prove metaphysical principles, there is an analogy and consistency.

\emph{Purusha} is also related to the Will, corresponding to the Hermetic principle\footnote{\url{https://gornahoor.net/?p=292}} ``Energy is concentrated Will''. Secular science seeks to effect material change by manipulation of energy.

Traditional science is based on the power of visualization and the strengthening of the Will.

\paragraph{Evolution}


Secular thought means two distinct things by ``evolution'', which on the surface seem to be inconsistent, yet really have the same presupposition. Biological theories assume evolution moves by random variation; nevertheless, there seems to be a ``direction'' to evolution, a ``tree of life'', that leads from less to more differentiation. However, were it truly random, there should be evidence of species evolving backwards. There should be examples of frogs returning to the sea, or isolated human groups becoming more ape-like.

In the social sense, ``to evolve'' means to become a conscious agent of some historic progression which is inevitable, not random, although some seem to evolve and some don't. The similarity of both positions is that ``the higher arises from the lower'', a metaphysical absurdity. Thus social evolutionists necessarily support an increase in animality and a decrease in differentiation. Metaphysically, this is really a regression to \emph{Prakriti} without \emph{Purusha}, a goal that is ultimately unrealizable and hence presents a natural barrier to the complete fulfillment of the idea of ``progress''.

In Tradition, ``to e-volve'' is understood in its etymological sense as an unfolding, that is, the unfolding in time of an eternal idea.

A Traditional science of evolution focuses on the understanding of metaphysical doctrines\footnote{\url{https://gornahoor.net/?p=823}} and how they manifest in time.

\paragraph{Unconscious}


Instead of leading to the freedom and emotional health of man, the probing of the unconscious by profane psychology has led to more powerful and effective methods to control his thoughts and behavior. Traditional science similarly recognizes a vital center in man which is likewise unconscious, because it is the subject of consciousness and cannot be its object. This center is transcendent and noumenal, or unknowable to profane science.

There are numerous traditional practices designed to train a man to live from that center. These may include developing the ability to concentrate or observing the flow of thoughts, feelings, desires without reacting to them. This center has been described in the West by mystics such as Eckhart or Boehme's ``The Nothing''. It is the source of freedom as Berdyaev writes:

\begin{quotex}
Freedom is not determined by God: it is part of the nothing out of which God created the world
\flright{\textit{The Destiny of Man}}
\end{quotex}


By concentrating and stilling the mind, an opening is prepared for the appearance of Being, or divine inspiration. This may also appear as a creative force to an artist.

The Traditional man is Free and lives an inspired and creative life.


\flrightit{Posted on 2011-01-04 by Cologero}

\begin{center}* * *\end{center}

\begin{footnotesize}
\begin{sffamily}
\texttt{Ismo Meinander on 2011-01-05 at 10:12 said:}

First of all, thank you Cologero and others for Gornahoor. It is truly a fountain source of inspiration and knowledge.

As it comes to this article and the notion of Quantum theory, have you read Wolfgang Smith's book ``The Quantum Enigma''? I think his ideas of vertical causality and of the ``form-bestowing agent'' are very interesting.


\hfill

\end{sffamily}
\end{footnotesize}

